\documentclass[11pt]{article}
\usepackage{microstyle}

\begin{document}

\lessonheader{6}{Chapter 3: Classical demand theory --- preferences, axioms, and utility}{}


\section{The classical approach}

We build demand behavior from scratch, by imposing fundamental axioms on consumer
preferences and seeing what these assumptions imply about behavior.


\section{The preference relation}

The primitive of the model is the preference relation $\succsim$. This is
hard-wired into the consumer and captures everything about his taste or
preference.

If $x,y\in X$ with $x\succsim y$, then $x$ is w.p.\ (weakly preferred) to $y$.
From $\succsim$:

\begin{flatlist}
  \item if $x\succsim y$ but not $y\succsim x$, then $x\succ y$;
  \item if $x\succsim y$ and $y\succsim x$, then the consumer is indifferent,
        $x\sim y$.
\end{flatlist}

\noindent
So we build $\succ$ and $\sim$ from $\succsim$; the latter is the primitive.


\section{The basic axioms: completeness and transitivity}

We start with axioms that imply that the consumer's preferences are relations in
the basic sense.

\begin{flatlist}
  \item \emph{Completeness:} for $x,y\in X$, either $x\succsim y$ or
        $y\succsim x$, or both.
  \item \emph{Transitivity:} if $x\succsim y$ and $y\succsim z$, then
        $x\succsim z$.
\end{flatlist}

This just tells us that the consumer is rational, but it does not tell us about
what his tastes are like. The remaining axioms have more economic meaning.


\section{Local non-satiation}

\runin{Local non-satiation.} For every $x\in X$ and any $\varepsilon>0$ there
exists $y\in X$ such that $\lVert x-y\rVert<\varepsilon$ and $y\succ x$.

So there is no bliss point: the consumer is never satiated, because he can always
get a better bundle close by. Two important implications:

\begin{flatlist}
  \item Indifference curves are thin --- kill the thick zone.

  \item The consumer, when he optimizes, spends his entire income: the budget
        constraint binds.
\end{flatlist}

\begin{figure}[ht]
\centering
\begin{tikzpicture}[scale=0.46]
  % ------------- panel (a): a thick indifference zone -------------
  \begin{scope}
    \ecaxes{4.2}{4.2}
    \node[ecnode,anchor=south west] at (0,4.25) {$x_2$};
    \node[ecnode,anchor=west]      at (4.35,0)   {$x_1$};

    \draw[budget] (0.5,3.3) .. controls (0.9,1.6) and (1.6,0.9) .. (3.4,0.6);
    \draw[budget] (1.2,3.7) .. controls (1.6,2.2) and (2.3,1.6) .. (3.9,1.3);
    \draw[line width=0.6pt] (0.72,2.55) -- (1.33,3.30) -- (1.00,1.68) -- (1.80,2.30)
          -- (1.40,1.25) -- (2.40,1.80) -- (2.25,0.90);
    \draw[red,fill=red!15] (1.62,1.55) circle (0.30);
    \node[ecnode,red] at (1.62,1.55) {$\varepsilon$};
  \end{scope}
  % ------------- panel (b): a thin indifference curve --------------
  \begin{scope}[xshift=6.4cm]
    \ecaxes{4.2}{4.2}
    \node[ecnode,anchor=south west] at (0,4.25) {$x_2$};
    \node[ecnode,anchor=west]      at (4.35,0)   {$x_1$};

    \draw[budget] (0.6,3.6) .. controls (1.05,1.75) and (1.70,1.10) .. (3.60,0.60);
    \draw[red,fill=red!15,fill opacity=0.6] (1.09,2.23) circle (0.34);
    \node[ecnode,red,anchor=west] at (1.43,2.23) {$\varepsilon$};
  \end{scope}
\end{tikzpicture}
\caption{The first implication of local non-satiation, drawn as in the original.
\emph{Left:} a thick indifference zone; a small $\varepsilon$-ball inside it
contains only bundles indifferent to its centre, so there is no strictly better
bundle nearby --- local non-satiation rules this out (``kill the thick zone'').
\emph{Right:} a thin indifference curve; every $\varepsilon$-ball around a point
on it reaches bundles above the curve.}
\end{figure}


\section{Why the budget constraint binds}

\runin{By contradiction.} Suppose the optimal bundle $x$ does not exhaust the
budget, $p\cdot x<w$. Then every bundle close enough to $x$ is still affordable
(there is an $\varepsilon>0$ such that $p\cdot y<w$ whenever
$\lVert y-x\rVert<\varepsilon$), and by local non-satiation one of them, $y$, has
$y\succ x$. So something better was affordable. This contradicts the fact that the
original bundle is optimal; hence $p\cdot x=w$.


\section{Monotonicity and convexity}

Axioms with more content regarding the structure of preferences.

\begin{flatlist}
  \item \emph{Monotonicity:} $\succsim$ is monotone if $x\gg y$ implies
        $x\succ y$.
  \item \emph{Strong monotonicity:} $\succsim$ is strongly monotone if
        $x\ge y$ and $x\ne y$, then $x\succ y$.
\end{flatlist}

Strong monotonicity implies monotonicity, and (on $X=\R^L_+$) monotonicity implies
local non-satiation, because the consumer would always like more of everything:
$x+\frac{\varepsilon}{2\sqrt{L}}(1,\dots,1)$ is within $\varepsilon$ of $x$ and strictly
preferred to it.

\begin{flatlist}
  \item \emph{Convexity:} $\succsim$ is convex if, for every $x$, $y\succsim x$ and
        $z\succsim x$ imply
        \[ \alpha y+(1-\alpha)z\succsim x\qquad\text{for all }\alpha\in[0,1]. \]
\end{flatlist}

Avr.\ (average) is always at least as good as extreme. Equivalently, for every $x$
the superior (upper contour) set is a convex set,
\[ \{\,y\in\R^L_+ : y\succsim x\,\}. \]
In a picture: indifference curves are bowed towards the origin (a diminishing
marginal rate of substitution).

\begin{figure}[ht]
\centering
\begin{tikzpicture}[scale=0.62]
  \fill[black!10] (0.6,3.6) .. controls (1.05,1.75) and (1.70,1.10) .. (3.60,0.60)
        -- (4.0,0.6) -- (4.0,4.0) -- (0.6,4.0) -- cycle;
  \ecaxes{4.4}{4.4}
  \node[ecnode,anchor=south west] at (0,4.45) {$x_2$};
  \node[ecnode,anchor=west]      at (4.5,0)   {$x_1$};
  \draw[budget] (0.6,3.6) .. controls (1.05,1.75) and (1.70,1.10) .. (3.60,0.60);

  \draw[densely dashed] (0.71,3.18) -- (2.83,0.84);
  \node[bundle] at (0.71,3.18) {};
  \node[ecnode,anchor=north east] at (0.66,3.12) {$y$};
  \node[bundle] at (2.83,0.84) {};
  \node[ecnode,anchor=north east] at (2.78,0.80) {$z$};
  \node[bundle] at (1.77,2.01) {};
  \node[ecnode,anchor=west] at (1.86,2.08) {$\alpha y+(1-\alpha)z$};
\end{tikzpicture}
\caption{Convexity, drawn as in the original: $y$ and $z$ lie on the
indifference curve through $x$, and their average $\alpha y+(1-\alpha)z$ lies in
the superior set of $x$ (hatched in the notes, shaded here), which is convex.}
\end{figure}


\section{Strict convexity}

\runin{Strict convexity.} $\succsim$ is strictly convex if, for every $x$,
$y\succsim x$, $z\succsim x$ and $y\ne z$ imply
\[ \alpha y+(1-\alpha)z\succ x \qquad \text{for all } \alpha\in(0,1). \]
Indifference curves then have no flat segments; with strictly convex (rational, continuous,
locally non-satiated) preferences the demand $x(p,w)$ is single-valued, a demand function.


\section{Utility functions}

The preference relation $\succsim$ is abstract; one may become used to thinking
about utility functions. A utility function is a map $u:X\to\R$ with
\[ x\succsim y \iff u(x)\ge u(y). \]
The utility function is just a made-up construct that is invented to represent
preferences in a different way. Utility functions are often more convenient to
work with, especially if we add assumptions such as differentiability.

Two facts follow directly from the definition. A utility function is \emph{ordinal}:
if $u$ represents $\succsim$, so does $f\circ u$ for any strictly increasing
$f:\R\to\R$, so only the ranking matters, not the numbers. And only a rational
$\succsim$ can be represented: $\ge$ on $\R$ is complete and transitive, so the
$\succsim$ it represents must be too.


\section{When a utility representation fails}

However, not all $\succsim$ can be represented by a utility function, even when
they are rational. The most famous example is lexicographic preference: put the
bundle with the most $x_2$, and compare $x_1$ only to break ties,
\[ x\succ y \iff x_2>y_2,\ \text{or}\ x_2=y_2\ \text{and}\ x_1>y_1 . \]
No two distinct bundles are indifferent, and there is no utility function for it.
It is not continuous: $x^n=(1,\,1+1/n)\succ y=(2,1)$ for every $n$, but in the limit
$x=(1,1)\prec y$.

To use a utility function, more assumptions are needed. At least, the extra
assumptions have a natural economic interpretation. We need continuity.


\section{Continuity}

\runin{Continuity.} $\succsim$ is continuous if for any sequence
$\{(x^n,y^n)\}_{n=1}^{\infty}$ with $x^n\succsim y^n$ for all $n$, and
$x=\lim_{n\to\infty}x^n$, $y=\lim_{n\to\infty}y^n$, it holds that
$x\succsim y$.

That is: no sudden preference reversal. Equivalently, for every $x$ the upper and
lower contour sets $\{y:y\succsim x\}$ and $\{y:x\succsim y\}$ are closed. The payoff
(Debreu): a rational and continuous $\succsim$ can be represented by a continuous
utility function. Lexicographic preferences fail exactly this continuity.

\begin{transcriptnotes}
  \item Page 1 of the original is a title page reading ``Lesson 6 : Chap.~3 ---
        classical demand Th.''; the abbreviation is read as \emph{theory}, and
        the page carries no other content. The subtitle previously read ``Part
        2''; it now follows the title page (Chapter 3).
  \item ``this hand written into the consumer'' has no verb in the original; it
        is read as \emph{This is} hard-wired into the consumer (``hand written''
        is taken as a mishearing of ``hard-wired'', the usual phrase).
  \item ``captures evenly about his taste or preference'' is read as
        \emph{captures everything} about his taste or preference --- the ending
        of the word is written as if it were ``-ly''.
  \item ``we start with axioms imply that the consumer's pref.\ are relations in
        basic sense'' is ungrammatical as written (the finite verb ``imply'' has
        no relative pronoun); it is read as \emph{axioms \textbf{that} imply
        \dots}. ``relations in basic sense'' is kept as written; in the notation
        of the course these are binary relations on $X$ in the ordinary
        set-theoretic sense.
  \item ``indifference curve are thin'' is corrected to \emph{curves}; ``kill the
        thick zone'' is written in red ink in the original.
  \item Local non-satiation is written with $y>x$ in the original; $y\succ x$ is
        used here, as elsewhere in these notes the strict relation is the hook
        symbol.
  \item \textbf{Corrected (earlier transcription):} page 4 draws two sketches
        under ``Two important implications'': on the left a thick indifference
        zone (a band hatched with a zig-zag) with a small red $\varepsilon$-ball
        inside it, on the right a thin indifference curve with a red
        $\varepsilon$-ball on it. Both red marks are $\varepsilon$-balls (the
        $\varepsilon$ is legible at high zoom), not a ``circled slash''. The
        figure previously had the panels swapped, a red slash, and invented
        labels $x$, $x'$; it now follows the original.
  \item ``By contradiction, the budget is not exhausted'' states the conclusion
        of a contradiction argument without spelling out the assumption; the
        sentence is read as \emph{suppose} the budget is not exhausted. The
        original also omits the verb in ``mean.\ that you can find something
        better'', read as \emph{meaning} that.
  \item ``rehanding'' $\to$ \emph{regarding}; ``AVR.'' is read as \emph{average}.
  \item \textbf{Corrected:} the convexity line was written with three conditions
        before the implication ($y\succsim x$, $y\succsim z$, $z\succsim x$) and no
        range for $\alpha$. The middle condition is redundant, and it is dropped;
        the definition is typeset in the standard form, $y\succsim x$ and
        $z\succsim x$ imply $\alpha y+(1-\alpha)z\succsim x$ for all
        $\alpha\in[0,1]$. ``Average is always better than extreme'' is softened to
        \emph{at least as good as}: convexity is a weak relation.
  \item The superior set is written $\{y\in\ \cdot\ \mid y\succsim x\}$ in the
        original; the set after $\in$ is an unclear glyph (like ``['' ), read as
        the consumption set $\R^L_+$. The convexity sketch marks three dots on
        the indifference curve region (two on the curve, their average inside the
        hatched set), reproduced as $y$, $z$ and $\alpha y+(1-\alpha)z$; a small
        pair of arrows drawn near the origin of that sketch is illegible and is
        omitted.
  \item ``stric convexity $\succsim$ is strict $\succsim$ is strictly convex''
        repeats the phrase; only \emph{strictly convex} is kept.
  \item \textbf{Corrected:} strict convexity is written ``if $y\succsim x$,
        $y\succsim z$, $z\succsim x$ implies $\alpha y+(1-\alpha)z\succ x$
        $\forall\alpha\in(0,1)$''. As with convexity, the middle condition is
        redundant and is dropped; the condition $y\ne z$ is missing and is added
        (without it $y=z=x$ would give $x\succ x$).
  \item \textbf{Added:} a precise proof that the budget constraint binds, the
        chain strong monotonicity $\Rightarrow$ monotonicity $\Rightarrow$ local
        non-satiation, the ordinal nature of utility and why representability
        requires rationality, the formal lexicographic order with a sequence showing
        it is not continuous, Debreu's representation result for continuous
        preferences, the equivalent closed-contour-set form of continuity, the
        remark that convex indifference curves are bowed toward the origin, and
        the remark after strict convexity (no flat segments; single-valued
        demand).
  \item ``is abstract, to may become use to thinking about utility functions'' is
        read as \emph{one may become used to thinking about utility functions};
        ``can be represented at a utility func.'' is read as \emph{by} a utility
        function.
  \item Lexicographic preference is written as ``put the bundle with the most
        $x_2$''. The order is normally described by comparing the first
        coordinate first; the notes' order (good 2 first) is kept, and the formal
        definition is written in that order.
  \item In the continuity definition the original labels the sequence
        $\{(x^n,y^n)\}_{N=1}^{\infty}$ and then writes ``$\forall N$'' before
        switching to lower-case $n$ in the limits. The index is written $n$
        throughout here.
  \item ``the extra assumptions has a nature econ interpretation'' is read as
        \emph{have a natural} economic interpretation.
  \item Page 9 of the original carries a single line, ``No sudden preference
        reversal.'', which closes the continuity section.
  \item Pages of the original, for tracing back: page 2 the classical approach
        and the preference relation; page 3 completeness, transitivity, and local
        non-satiation; page 4 the two implications with the thick-zone sketch;
        page 5 the binding budget constraint and monotonicity; page 6 strong
        monotonicity, convexity and the superior set; page 7 strict convexity and
        the utility function; page 8 lexicographic preference and continuity;
        page 9 the closing line.
\end{transcriptnotes}

\end{document}
